\section{Tata Bahasa sebagai Sistem Penulisan Ulang}
\label{sec:ch02-grammar}

\subsection{Definisi Tata Bahasa}

\begin{definisi}[Tata bahasa]
Tata bahasa \(G\) didefinisikan sebagai empat-tupel
\[
G = (V, T, P, S),
\]
dengan \(V\) himpunan berhingga simbol nonterminal, \(T\) himpunan
berhingga simbol terminal, \(P\) himpunan aturan produksi, dan
\(S \in V\) simbol awal. Himpunan \(V\) dan \(T\) saling lepas.
\end{definisi}

Aturan produksi berbentuk \(\alpha \to \beta\), dengan \(\alpha\)
memuat paling sedikit satu nonterminal. Jika \(\alpha\) hanya berupa
satu nonterminal, aturan disebut \emph{context-free}. Bentuk khusus
\(A \to \alpha\) inilah yang menjadi fokus Bab VIII dan seterusnya.

Notasi BNF (\emph{Backus-Naur Form}) menulis aturan
\(A \to \alpha_1 \mid \alpha_2 \mid \cdots\) sebagai satu aturan
dengan alternatif. Sebagai contoh, tata bahasa aritmetika sederhana
dapat ditulis:
\[
\begin{aligned}
E &\to E + T \mid T,\\
T &\to T * F \mid F,\\
F &\to (E) \mid \code{id}.
\end{aligned}
\]

\subsection{Derivasi dan Bentuk Kalimat}

\begin{definisi}[Derivasi]
Jika \(A \to \gamma\) adalah aturan produksi dan
\(\alpha A \beta\) adalah string simbol, maka
\(\alpha A \beta \Rightarrow \alpha \gamma \beta\) adalah satu
langkah derivasi. Notasi \(\Rightarrow^{*}\) menyatakan nol atau lebih
langkah derivasi. String yang diperoleh dari \(S\) melalui derivasi
disebut bentuk kalimat; bentuk kalimat yang seluruhnya terminal
disebut string dalam bahasa.
\end{definisi}

Bahasa yang dibangkitkan oleh \(G\) adalah
\[
L(G) = \{w \in T^{*} \mid S \Rightarrow^{*} w\}.
\]

Derivasi \emph{leftmost} selalu mengganti nonterminal paling kiri;
derivasi \emph{rightmost} mengganti nonterminal paling kanan. Untuk
tata bahasa yang tidak ambigu, kedua derivasi menghasilkan pohon urai
yang sama.

\begin{contoh}[Derivasi leftmost]
Untuk string \(\code{id} + \code{id} * \code{id}\), derivasi leftmost
pada tata bahasa aritmetika adalah:
\[
\begin{aligned}
E &\Rightarrow E + T\\
  &\Rightarrow T + T\\
  &\Rightarrow F + T\\
  &\Rightarrow \code{id} + T\\
  &\Rightarrow \code{id} + T * F\\
  &\Rightarrow \code{id} + F * F\\
  &\Rightarrow \code{id} + \code{id} * F\\
  &\Rightarrow \code{id} + \code{id} * \code{id}.
\end{aligned}
\]
\end{contoh}

\subsection{Pohon Urai}

Derivasi dapat direpresentasikan sebagai pohon urai (\eng{parse tree}).
Akar pohon adalah simbol awal \(S\). Setiap simpul dalam adalah
nonterminal, dan anak-anaknya adalah simbol pada ruas kanan aturan
produksi yang dipakai. Daun-daun pohon, dibaca dari kiri ke kanan,
membentuk string yang diturunkan.

\begin{figure}[H]
\centering
\begin{forest}
parsetree
[E
  [E [T [F [\code{id}]]]]
  [+]
  [T
    [T [F [\code{id}]]]
    [*]
    [F [\code{id}]]
  ]
]
\end{forest}
\caption{Pohon urai untuk \(\code{id} + \code{id} * \code{id}\)}
\label{fig:ch02-pohon-urai}
\end{figure}

Pohon urai menunjukkan struktur hierarkis ekspresi: perkalian lebih
dalam daripada penjumlahan, sehingga dikerjakan lebih dahulu. Struktur
inilah yang nantinya diringkas menjadi \emph{Abstract Syntax Tree}
(AST) pada Bab XVI.

\subsection{Klasifikasi Tata Bahasa}

Tata bahasa dapat diklasifikasikan menurut bentuk aturan produksinya.
Tata bahasa regular memiliki aturan berbentuk \(A \to aB\) atau
\(A \to a\), tata bahasa bebas konteks memiliki aturan
\(A \to \alpha\), tata bahasa sensitif konteks memiliki aturan
\(\alpha A \beta \to \alpha \gamma \beta\) dengan
\(\gamma \neq \varepsilon\), dan tata bahasa tanpa batas tidak
membatasi bentuk aturan. Klasifikasi ini bersesuaian dengan hierarki
Chomsky yang dibahas pada bagian berikutnya.
